% ============== Section 4: Ultra-massive MIMO & beamforming (S21-S29) =====
\section{Ultra-Massive MIMO \& Beamforming}

% --- S21 -----------------------------------------------------------------
\begin{frame}{From massive to ultra-massive}
  \begin{itemize}\itemsep2pt
    \item 5G today: 64\,TRX at 3.5\,GHz \;$\to$\; 6G: \alert{512--4096
          elements} at FR3 / mmWave / sub-THz
    \item Same rooftop aperture at higher $f$ $\Rightarrow$ more elements,
          narrower beams, more EIRP
    \item New physics enters: \alert{near-field propagation}, spatial
          non-stationarity across large arrays
    \item The bill: RF chains, ADCs, calibration, fronthaul --- DSP must
          economize
  \end{itemize}
  \note{[S21, 1 min | clock 0:26:30, after Demo 1] Re-entry after the demo:
  ``From time-frequency games to space.'' The scaling is driven by aperture
  reuse: same panel, quadratically more elements as wavelength shrinks. The
  section's plot: every gain has a hardware price, and the DSP's job is to
  refuse to pay full retail.}
\end{frame}

% --- S22 -----------------------------------------------------------------
\begin{frame}{Near-field vs far-field}
  \begin{columns}[T]
    \column{0.52\textwidth}
    \begin{equation*}
      d_R=\frac{2D^{2}}{\lambda}
      \qquad
      \begin{cases}
        d>d_R: & \text{planar, }\mathbf a(\theta)\\[2pt]
        d<d_R: & \text{spherical, }\mathbf a(\theta,r)
      \end{cases}
    \end{equation*}
    \begin{itemize}\itemsep1pt
      \item $D=0.5$\,m at 30\,GHz: $d_R=50$\,m --- users are
            \emph{inside} the near field
      \item Near field $\Rightarrow$ \alert{beam focusing} (range-selective!),
            not just steering
      \item Wavefront curvature carries range information --- gift for
            localization
    \end{itemize}
    \column{0.46\textwidth}
    \centering
    \begin{tikzpicture}[scale=0.75]
      \draw[very thick, accent] (0,0) -- (0,2.4);
      \foreach \y in {0.2,0.6,...,2.3} \fill[accent] (0,\y) circle (1.1pt);
      \node[lab, below] at (0,-0.1) {array, aperture $D$};
      \foreach \r in {1.0,1.6,2.2}
        \draw[cblue, thick] (0,1.2) ++(-62:\r) arc (-62:62:\r);
      \node[lab, cblue] at (1.5,2.6) {spherical};
      \foreach \x in {3.6,4.1,4.6}
        \draw[corange, thick] (\x,0.1) -- (\x,2.3);
      \node[lab, corange] at (4.1,2.6) {planar};
      \draw[arr, gray] (0.1,0.05) -- node[lab, below]{$d_R=2D^2/\lambda$} (3.5,0.05);
    \end{tikzpicture}
  \end{columns}
  \note{[S22, 1.5 min | clock 0:27:30] Everyone learned $d_R$ once and forgot
  it because it used to be centimeters. Recompute it live: half-meter array,
   30 GHz, 50 meters --- the cell IS the near field. Beam focusing means two
  users on the same bearing at different ranges can be separated. That is a
  genuinely new spatial degree of freedom.}
\end{frame}

% --- S23 -----------------------------------------------------------------
\begin{frame}{Hybrid analog--digital architectures}
  \centering
  \resizebox{0.9\textwidth}{!}{% Hybrid analog-digital beamforming architecture (slide 23) -- native TikZ
\begin{tikzpicture}[node distance=3mm and 7mm]
  % ---------------- fully digital (left) ----------------
  \node[lab] (s1) {$N_s$ streams};
  \node[blkb, right=of s1, minimum height=16mm] (bb1) {Digital\\precoder};
  \node[blk, right=of bb1, yshift=7mm, minimum width=12mm] (rf1) {RF chain};
  \node[blk, right=of bb1, yshift=-7mm, minimum width=12mm] (rf2) {RF chain};
  \node[lab] at ($(rf1)!0.5!(rf2)$) {$\vdots\; N_t$};
  \foreach \i/\n in {rf1/a1, rf2/a2}{
    \node[right=4mm of \i] (\n) {\antsym};}
  \draw[arr] (s1) -- (bb1);
  \draw[arr] (bb1.east|-rf1) -- (rf1); \draw[arr] (bb1.east|-rf2) -- (rf2);
  \draw[arr] (rf1) -- (a1); \draw[arr] (rf2) -- (a2);
  \node[lab, align=center] at ($(bb1)!0.55!(a1)+(0,14mm)$)
    {\textbf{Fully digital}: one RF chain \emph{per antenna}\\
     $N_t$ ADC/DAC pairs -- power \& cost prohibitive at high bands};

  % ---------------- hybrid (right) ----------------
  \begin{scope}[xshift=72mm]
  \node[lab] (s2) {$N_s$ streams};
  \node[blkb, right=of s2, minimum height=16mm] (bb2) {$F_{\mathrm{BB}}$\\ \scriptsize digital};
  \node[blk, right=of bb2, yshift=7mm, minimum width=12mm] (rf3) {RF chain};
  \node[blk, right=of bb2, yshift=-7mm, minimum width=12mm] (rf4) {RF chain};
  \node[lab] at ($(rf3)!0.5!(rf4)$) {$N_{\mathrm{RF}}\!\ll\! N_t$};
  \node[blk, right=17mm of bb2, minimum height=16mm, minimum width=13mm] (ps)
    {$F_{\mathrm{RF}}$\\ \scriptsize phase\\[-3pt]\scriptsize shifters};
  \node[right=4mm of ps, yshift=7mm] (a3) {\antsym};
  \node[right=4mm of ps] (a4) {\antsym};
  \node[right=4mm of ps, yshift=-7mm] (a5) {\antsym};
  \node[lab, right=1mm of a4, xshift=2mm] {$N_t$};
  \draw[arr] (s2) -- (bb2);
  \draw[arr] (bb2.east|-rf3) -- (rf3); \draw[arr] (bb2.east|-rf4) -- (rf4);
  \draw[arr] (rf3) -- (ps.west|-rf3); \draw[arr] (rf4) -- (ps.west|-rf4);
  \draw[arr] (ps.east|-a3) -- (a3); \draw[arr] (ps.east|-a4) -- (a4);
  \draw[arr] (ps.east|-a5) -- (a5);
  \node[lab, align=center] at ($(ps)+(-8mm,17.5mm)$)
    {\textbf{Hybrid}: few RF chains + analog phase network\\
     $y = W_{\mathrm{BB}}^{H}W_{\mathrm{RF}}^{H}
          \left(H F_{\mathrm{RF}}F_{\mathrm{BB}}\,s + n\right)$};
  \end{scope}
\end{tikzpicture}
}
  \begin{itemize}\itemsep1pt
    \item $F_{\mathrm{RF}}$: phase-only, $|[F_{\mathrm{RF}}]_{ij}|=1/\sqrt{N_t}$;
          wideband arrays add \alert{true-time-delay} units against beam squint
  \end{itemize}
  \note{[S23, 1.5 min | clock 0:29] Left: the fully digital ideal --- an
  ADC/DAC pair per element, unaffordable at 1024 elements. Right: the hybrid
  compromise --- a thin digital layer, a cheap analog phase network. The
  constraint that makes the math interesting: RF entries are unit-modulus.
  Mention subarray (partially connected) variant trades gain for phase-shifter
  count.}
\end{frame}

% --- S24 -----------------------------------------------------------------
\begin{frame}{Designing the hybrid precoder: sparsity to the rescue}
  \begin{equation*}
    \min_{F_{\mathrm{RF}},F_{\mathrm{BB}}}
    \left\|F_{\mathrm{opt}}-F_{\mathrm{RF}}F_{\mathrm{BB}}\right\|_F
    \quad\text{s.t.}\quad
    F_{\mathrm{RF}}\in\{\text{array response vectors}\}
  \end{equation*}
  \begin{columns}[T]
    \column{0.5\textwidth}
    \begin{itemize}\itemsep1pt
      \item mmWave channels are sparse in angle $\Rightarrow$ precoding
            $\approx$ \alert{sparse reconstruction} $\Rightarrow$ OMP
            (El Ayach et al.\ 2014)
      \item $N_{\mathrm{RF}}\approx N_s$ chains already $\approx$
            near-digital performance
    \end{itemize}
    \column{0.48\textwidth}
    \centering
    \includegraphics[width=\textwidth]{fig_hybrid_se.png}
  \end{columns}
  \note{[S24, 1 min | clock 0:30:30] The curves are our own simulation (same
  code as Demo 2): 64x16, 6 paths, 4 streams --- the 4-RF-chain OMP hybrid
  sits within a fraction of a bit of fully digital. The deep reason: the
  channel only has ~6 directions worth talking to; 4 well-chosen beams capture
  almost everything.}
\end{frame}

% --- S25 -----------------------------------------------------------------
\begin{frame}{Beam management: finding the needle}
  \begin{itemize}\itemsep2pt
    \item 5G NR baseline: SSB sweep $\to$ CSI-RS refinement $\to$ tracking
          and failure recovery
    \item Overhead grows with codebook size $\Rightarrow$ hierarchical search,
          ML-assisted prediction for 6G
    \item Sub-THz pencil beams: mobility and micro-blockage stress the loop
          hard
    \item Research directions: side-information beams (position, vision,
          radar), \alert{near-field codebooks} (angle $\times$ range)
  \end{itemize}
  \note{[S25, 1 min | clock 0:31:30] Beam management is where beautiful array
  theory meets protocol reality: every training symbol is a symbol not
  carrying data --- recall mmWave Lesson 2. The 6G twist: near-field codebooks
  must sample angle AND range, squaring the codebook unless structure is
  exploited. Natural hook for the AI section later.}
\end{frame}

% --- S26 -----------------------------------------------------------------
\begin{frame}{Cell-free massive MIMO}
  \begin{columns}[T]
    \column{0.55\textwidth}
    \begin{itemize}\itemsep2pt
      \item Many distributed low-cost APs, coherent joint processing, no cell
            boundaries
      \item Uniform SNR replaces cell-edge misery; macro-diversity beats
            blockage
      \item The catches: fronthaul, synchronization, scalable processing
            (dynamic cooperation clusters)
      \item 6G angle: natural platform for \alert{distributed sensing} and RIS
            integration
    \end{itemize}
    \column{0.42\textwidth}
    \centering
    \begin{tikzpicture}[scale=0.72]
      % cellular
      \foreach \x/\y in {0/0, 1.6/0, 0.8/1.35}
        \node[regular polygon, regular polygon sides=6, draw=gray,
              minimum size=15mm, inner sep=0] at (\x,\y) {};
      \node[lab] at (0.8,-0.85) {cells + towers};
      % cell-free
      \begin{scope}[xshift=38mm]
        \node[blkb, minimum width=9mm] (cpu) at (0.8,-0.1) {CPU};
        \foreach \x/\y in {0/0.5, 0.5/1.5, 1.3/0.9, 1.9/1.6, 2.1/0.3, 1.1/2.0}{
          \fill[accent] (\x,\y) circle (2.2pt);
          \draw[gray!60, thin] (cpu) -- (\x,\y);}
        \node[lab] at (1.1,-0.85) {APs + fronthaul};
      \end{scope}
    \end{tikzpicture}
  \end{columns}
  \note{[S26, 1 min | clock 0:32:30] One-line pitch: instead of a few big
  towers each owning users, many small APs all serve everyone. The
  90-percentile user is the winner. Physical-layer research lives in the
  ``catches'' bullet --- especially scalable combining and phase sync over
  fronthaul.}
\end{frame}

% --- S27 -----------------------------------------------------------------
\begin{frame}{Channel estimation by compressed sensing}
  \begin{equation*}
    \mathbf y=\Phi\,\mathbf h+\mathbf n,\qquad
    \|\mathbf h\|_0=L\ll\dim(\mathbf h),\qquad
    H\approx A_R\,\mathrm{diag}(\mathbf g)\,A_T^{H}
  \end{equation*}
  \begin{itemize}\itemsep2pt
    \item Angular-domain sparsity ($L$ paths) $\Rightarrow$ recovery from
          $\mathcal O(L\log N)$ pilots via OMP / AMP
    \item Pilot savings are what make 1024-element arrays \emph{trainable at
          all}
    \item Caveats: grid mismatch (off-grid/atomic-norm methods), near-field
          dictionaries grow with angle$\times$range
    \item Same math returns in Section 8 for activity detection --- learn once,
          bill twice
  \end{itemize}
  \note{[S27, 1 min | clock 0:33:30] Connect to S9: the physics gave us
  sparsity, CS cashes it in. Order-of-magnitude example: 64 paths worth of
  unknowns in a 1024x1 vector needs a few hundred, not a few thousand,
  measurements. Flag the two honest caveats; off-grid is a good student
  project generator.}
\end{frame}

% --- S28 -----------------------------------------------------------------
\begin{frame}{Pilot contamination}
  \begin{equation*}
    \mathrm{SINR}_j\;\xrightarrow{\;M\to\infty\;}\;
    \frac{\beta_{jj}^{2}}{\sum_{l\neq j}\beta_{jl}^{2}}
    \qquad\text{(same-pilot cells }l\text{)}
  \end{equation*}
  \begin{itemize}\itemsep2pt
    \item Finite coherence block $\Rightarrow$ pilot reuse $\Rightarrow$
          correlated estimates $\Rightarrow$ coherent interference that does
          \alert{not} average out with more antennas
    \item Modern view (Bj\"ornson et al.): with multi-cell MMSE processing and
          spatial correlation, capacity grows without bound ---
          contamination is a \emph{processing artifact}, not a law of nature
    \item Practical mitigations: smart pilot assignment, correlation-aware
          estimation, cell-free operation
  \end{itemize}
  \note{[S28, 1 min | clock 0:34:30] Tell it as a two-act story: 2010 --- the
  asymptotic ceiling, most-cited limit in massive MIMO; 2018 --- the ceiling
  dissolves under MMSE processing with unequal spatial correlations. Great
  case study in how ``fundamental limits'' sometimes encode assumptions, and
  how better DSP quietly removes them.}
\end{frame}

% --- S29 -----------------------------------------------------------------
\begin{frame}{Array frontier: open challenges}
  \begin{itemize}\itemsep2pt
    \item Wideband beam \alert{squint} at $10\%+$ fractional bandwidth:
          TTD-augmented hybrids
    \item Near-field CSI: dimensionality, codebooks, tracking
    \item Spatial non-stationarity: different subarrays see different channels
    \item The honest metric: \alert{energy per beamformed bit}, not array gain
  \end{itemize}
  \demoslot{2 (3.5 min)}{\textbf{64-element array, live}: steer the beam with a
    slider; digital vs OMP-hybrid spectral efficiency.
    \hfill\src{matlab/demo2\_hybrid\_beamforming}}
  \note{[S29 + DEMO 2, 1+3.5 min | clock 0:35:30] Quick pass over the open
  list, then SWITCH TO MATLAB. Demo 2 script: move the slider, the beam
  follows; grating-lobe moment when steering beyond 60 degrees is a nice
  ``physics is watching'' beat. Then the SE bars: 4 chains vs 64. If MATLAB
  fails: appendix A2. Demo done by 0:40.}
\end{frame}
